\begin{statement}{E.38}{Proved}{The complete double-cover Ferrari refinement retains every fibre}
Put $L=\operatorname{SL}_2(\F_5)$ and retain all epimorphism
kernels of the free Ferrari group onto $L$.  The resulting
universal free image is
\[
             \mathcal L=L^{19}\times\F_2^2.
 \tag{E.38.1}
\]
The literal Ferrari action gives a finite quotient
$B_4\twoheadrightarrow T_L$.  Inducing its six evaluation
maps through the six cosets of $F\lhd B$ gives an actual
finite quotient $\rho_L:B\twoheadrightarrow G_L$.
Its common refinements with E.24 have the exact descriptions
\[
 D_{23}\mathbin{\times_{H_L}}G_L,
 \qquad
 T_{23}\mathbin{\times_{S_L}}T_L,
 \qquad |H_L|=1728,
 \quad |S_L|=96.
 \tag{E.38.2}
\]
Here $H_L$ is the specified quotient of $G_6$, and $S_L$
is the maximal solvable quotient of $T_L$ with its
specified map from the binary degree-two Ferrari target.

Every one of the $864$ finite points of E.24 has a correction
in $N'$ preserving its complete $D_{23}$ coordinate and
satisfying the source hexagons, the full pentagon, the braid
relations and the full-twist identities in these common
refinements.  The source and four-strand maps are surjective.
This holds for each compatible profinite unit multiplier:
its residue modulo $144$ selects the corresponding E.24
branch.  Thus every retained CM fibre, including every extra
fibre, survives this central-cover test.
\end{statement}
\begin{proof}
We give the universal comparison before the correction
calculation.  The general construction by generating-pair
orbits is described in \cite[\S2.1, Lemma 2.1]{DE_Guillot2015};
the following identifies its actual image in this case.

\paragraph{The universal free quotient and its action.}
The center of $L$ is $Z=\{I,-I\}$ and $L/Z=A_5$.
All $120$ matrices, their products, their center, and the
equality $L'=L$ are checked directly.  The action of $L/Z$
on its five Sylow $2$-subgroups identifies it with the
specific $A_5$ model of E.29.  Conjugation by
$\operatorname{PGL}_2(\F_5)$ realizes all $120$
automorphisms of that model.  These are also all
automorphisms of $L$: an automorphism trivial on $L/Z$
has the form $g\mapsto g\epsilon(g)$ with
$\epsilon:L\to Z$ a homomorphism, which is trivial
because $L$ is perfect.

Each of the $19$ epimorphism kernels onto $A_5$ has
four lifts onto $L$, obtained by the two independent
signs of its generator images.  Every lifted pair
generates $L$.  Indeed, its image $J$ has $JZ=L$,
and $L=[JZ,JZ]=[J,J]\subset J$.
The four kernels are distinct, since an automorphism
identifying two lifts would be trivial on $A_5$ and
hence trivial on $L$.  Conversely every epimorphism
onto $L$ has one of these reductions.  There are
therefore exactly $76$ kernels.

Choose one lift $(g_j,h_j)$ over each of the $19$
$A_5$ pairs.  The joint image $J$ in $L^{19}$
surjects onto $A_5^{19}$ by E.29, so
$JZ^{19}=L^{19}$.  Perfectness again gives
$J=L^{19}$.  Adjoining the two free abelianization
coordinates modulo $2$ gives the full product
\textup{(E.38.1)}: a common quotient of a perfect
group and an abelian group is trivial.  All $76$
evaluations are recovered by the actual formula
\[
 \theta_{j,e_1,e_2}(w)
   =(-I)^{e_1\operatorname{ab}_1(w)+
               e_2\operatorname{ab}_2(w)}\theta_j(w).
 \tag{E.38.3}
\]
Conversely the chosen $19$ values recover $L^{19}$,
and ratios with the sign-twisted values recover
both parity coordinates.  Thus this formula gives
an isomorphism from $\mathcal L$ onto the actual
joint image in $L^{76}$.

Write a Nielsen automorphism as $(p,c,D,M)$, acting by
\[
 (g,t)\longmapsto
 \bigl(((-I)^{D_jt}c_jg_{p(j)}c_j^{-1})_j,Mt\bigr),
 \quad
 D\in\operatorname{Mat}_{19\times2}(\F_2),\quad
 M\in\operatorname{GL}_2(\F_2).
 \tag{E.38.4}
\]
Here $c_j\in\operatorname{PGL}_2(\F_5)$ acts by the
actual matrix conjugation.  For the substitutions
$a:(x,y)\mapsto(x,yx)$ and
$b:(x,y)\mapsto(xy^{-1},y)$, the permutation and
the $c_j$ values are those of E.29.  Their two
lifted generator images determine every entry of
$D$, while the matrices $M$ are respectively
$\left[\begin{smallmatrix}1&1\\0&1\end{smallmatrix}\right]$
and
$\left[\begin{smallmatrix}1&0\\1&1\end{smallmatrix}\right]$.
Composition is the literal rule
\[
 (p,c,D,M)(q,d,E,R)
   =\bigl(q\circ p,(c_jd_{p(j)})_j,
                  DR+E\circ p,MR\bigr).
 \tag{E.38.5}
\]
For the semidirect product, the free coordinate
multiplies as
$k_j(-I)^{D_ju}c_jl_{p(j)}c_j^{-1}$, and the
parity coordinate as $t+Mu$.
These formulas retain the possible parity-to-center
cross terms.  The three braid generators are
$(1,a),(1,b),(\xi_1,a)$ as in E.29; the literal
Ferrari identities and every braid relation are
checked in this model.

The two exact projections are to $T_A$ and to
$\operatorname{AGL}_2(\F_2)=S_4$, the latter
given by $(M,t)$.  Their joint image is
$T_A\times_{C_2}S_4$: E.29 identifies the maximal
solvable quotient of $T_A$ with braid parity,
and Goursat's lemma applies to these specified
generator maps.  Above its identity, the $19$
free central signs and $38$ cross-term signs
add in $\F_2^{57}$.  The actual kernel has
dimension $55$, consisting of the $19$ free
signs and the computed $36$-dimensional Nielsen
kernel.  In particular
\[
             |T_L|=2^{117}3^{36}5^{31}7^2.
 \tag{E.38.6}
\]

\paragraph{The complete induced source and correction groups.}
For the six literal evaluation maps
$\theta_i:F\to T_L$, use the same section of
$B/F$ as in E.29.  The induced component is
\[
 (q,i)\longmapsto\theta_i(b_qwb_{qw}^{-1}),
 \tag{E.38.7}
\]
together with the permutation $q\mapsto qw$.
This defines the homomorphism $\rho_L$ of discrete
groups and its continuous extension from $\widehat B$.
Its kernel is normal in $B$ and lies in $F$.
The identity-coset coordinate on $\rho_L(F)$ is
the specified map to each evaluation target.
Projection to E.29 gives a surjection
$\pi:G_L\to G_A$ with elementary abelian kernel
of dimension $44$.

Set $P=\rho_L(N)$, $C=\rho_L(N')$ and
$E=\rho_L(F')$.  The exact image calculations give
\[
 \begin{array}{c|c|c}
 \text{group}&\text{order}&
       \dim_{\F_2}\ker(\pi|_{\text{group}})\\\hline
 G_L&2^{82}3^{26}5^{21}7^2&44\\
 P&2^{76}3^{24}5^{21}7^2&41\\
 C&2^{76}3^{23}5^{19}7^2&41\\
 C'&2^{76}3^{23}5^{19}7^2&41\\
 E&2^{79}3^{23}5^{19}7^2&42
 \end{array}
 \tag{E.38.8}
\]
We describe completeness of the computation.
Above a faithful permutation quotient, each
stabilizer-chain Schreier generator is rewritten
through the next transversals.  Its discrepancy
is evaluated in the actual elementary kernel.
All such discrepancies, with their normal
conjugates, give the kernel of the lifted group:
successive Schreier rewriting reduces every
identity-base word to their normal closure.
Every retained generator is an actual relator,
recorded in a word DAG which is independently
substituted at the end of the program.

For $P$, all $97$ original $N$ generators are
retained, including the $91$ differences above
the six selected base generators.  To compute
the derived group of $\langle g_i,V\rangle$,
where $V$ is the complete elementary kernel,
we retain every $[g_i,g_j]$, every $[v,g_i]$
for a basis of $V$, every conjugation residual
of the selected base-derived generators, and
the complete lifted base presentation.
Their normal closure is exactly the derived
group: the resulting quotient is abelian,
and every retained word is a commutator or
its conjugate.  Applying this construction
twice proves $C'=C$.  Module commutators alone
have rank $34$ in its $41$-dimensional kernel;
the remaining seven directions are retained
by the full second-derived calculation.
Exhaustive calculations in $A_4$ and in the
order-eight Heisenberg group independently
check the lifted-presentation procedure,
including an extra lift and its normal closure.

The equality of the two $41$-dimensional
kernels in \textup{(E.38.8)} proves
\[
 \pi^{-1}(C_A)\cap P=C,
 \qquad P/C\simeq P_A/C_A.
 \tag{E.38.9}
\]
Thus E.29's complete correction calculation
transfers through this displayed isomorphism:
with $K=\rho_L(K_6)$,
\[
 |P/C|=75,\quad |K/C|=25,\quad
 \rho_L(K_6\cap F')=C,\quad |G_L/K|=1728.
 \tag{E.38.10}
\]
For the third equality, the projection of
$K_6\cap F'$ lies in $C_A$, hence its new
image lies in $C$ by \textup{(E.38.9)};
the reverse inclusion follows from
$N'\subset K_6\cap F'$.

\paragraph{The signed modular comparison.}
Use the literal map
\[
 q_{12}:B\longrightarrow Q=
       \operatorname{SL}_2(\mathbb Z/12)/\{\pm I\},
 \qquad
 s\longmapsto\begin{bmatrix}0&1\\-1&0\end{bmatrix},
 \quad r\longmapsto\begin{bmatrix}0&-1\\1&-1\end{bmatrix}.
 \tag{E.38.11}
\]
All $97$ defining $N$ words have identity image,
and $[G_L:P]=576=|Q|=[B:N]$.
Since $N$ is contained in the preimage of $P$,
the equality of indices proves that this
preimage is exactly $N$.  Thus the displayed
map descends to $G_L$, with kernel $P$.

In $\operatorname{SL}_2(\mathbb Z/12)$ the actual
free generators are
$x=\left[\begin{smallmatrix}1&2\\0&1\end{smallmatrix}\right]$
and
$y=\left[\begin{smallmatrix}1&0\\-2&1\end{smallmatrix}\right]$.
They generate a group of order $96$ whose
derived subgroup has order eight.  Every
matrix of that derived subgroup is $I$
modulo four.  Consequently its projective
image has a unique representative with that
property, and its reduction modulo three
is the quaternion subgroup $Q_8$ of
$\operatorname{SL}_2(\F_3)$.

There is also an independent proof of the
last row's quotient comparison in
\textup{(E.38.8)}.  E.29 gives
$E_A\cap P_A=C_A$.  Combining this equality
with \textup{(E.38.9)} shows $E\cap P=C$.
The explicit modular calculation therefore gives
\[
                 E/C\xrightarrow{\sim}q_{12}(E)=Q_8.
 \tag{E.38.12}
\]
The direct full $F'$ computation agrees, with
kernel dimension $42$.

The $24$ retained E.17 words have the following
three signed values, including their actual
$\operatorname{SL}_2$ lifts:
\[
 M_0=I,\qquad
 M_1=\begin{bmatrix}1&4\\4&5\end{bmatrix},\qquad
 M_2=\begin{bmatrix}1&8\\8&5\end{bmatrix}
                      \pmod{12}.
 \tag{E.38.13}
\]
Their complete index table is (E.29.13).
The modular program evaluates every one of
the $24$ old correction words and all six
ternary and seven binary fibre-basis words.
Every value is identity in $Q$; it also
checks the complete base expression with
its actual integral cusp gauge.  Hence
all $864$ reconstructed words have exactly
the indexed value in \textup{(E.38.13)}.
This proves the whole-fibre assertion
through the word reconstruction formula.

\paragraph{Arithmetic selection and a common correction.}
Fix a compatible profinite unit $\lambda$.
Its old residue is $1$ or $73$ modulo $144$,
so $\lambda\equiv1\pmod{12}$.  Cyclotomic
surjectivity supplies an actual Galois
element with character $\lambda$.
As proved in E.29,
$\mathbb Q(\zeta_3,\sqrt[3]{2})\cap
 \mathbb Q(\mu_\infty)=\mathbb Q(\zeta_3)$.
Multiplication by an element of the
cyclotomic kernel therefore makes its
Kummer value $\rho_2\bmod3$ any prescribed
$\kappa$.  Nakamura--Schneps
\cite[Theorem 2.5]{BCNakamuraSchneps2000},
applied to this actual Galois element,
gives exactly $q_{12}(f_\sigma)=M_\kappa$.
The sign in that theorem is positive
because $\lambda\equiv1\pmod4$.
No extension of that arithmetic matrix
formula to arbitrary GT elements is used.

For a retained word $f_0\in F'$ choose
the matching $\kappa$.  By
\textup{(E.38.12)},
\[
             c=\rho_L(f_0)^{-1}\rho_L(f_\sigma)\in C.
 \tag{E.38.14}
\]
The image of $N'$ in $D_{23}$ is solvable.
Its joint image with the perfect group $C$
is the full product by Goursat's lemma,
since a perfect solvable quotient is trivial.
There is therefore a discrete $k\in N'$
with these two images equal to $c$ and $1$.
The word $f=f_0k$ has exactly its old
$D_{23}$ coordinate and exactly the
$G_L$ coordinate of $f_\sigma$.
It remains in $F'$.
All $36$ induced coordinates match, so
the six target evaluations, both source
hexagons, the full pentagon and the braid
and twist equations agree with the actual
Galois action in the new component.
The old component retains its E.24 equations.
The braid convention is the exact
Nakamura--Schneps (1.1) comparison proved
at the end of E.29, including the sixth pair.

The multiplier comparison includes every
power in these assertions.  In $D_{23}$
the cusp $t$ has order $144$ and the three
conjugate peripheral generators $x,y,z$
have order $72$.  Thus the residue of
$\lambda$ modulo $144$ fixes both the
required cusp power and every power with
$m=(\lambda-1)/2$.  The old four-strand
generators have orders dividing eight
and nine, and their full-twist periods
were retained in E.24.  The outer
permutation quotient is the same $S_3$.
In the new component we use the very same
profinite $\lambda$ as the actual Galois
element, so no further power replacement
is needed.  Explicit representatives for
the two finite branches are $1$ and the
unit $\eta_{73}$ whose $p$-adic component
is $73$ for $p=2,3,5,7$ and $1$ otherwise.
All the finite groups here have only those
prime divisors, so $\eta_{73}$ and
$(\eta_{73}-1)/2$ give exactly the integer
powers $73$ and $36$ in every component.

\paragraph{Joint generation.}
Set $H_L=G_L/K$.  Its map from $B$ factors
through $G_6$, and hence through $D_{23}$.
Every common quotient of $D_{23}$ and
$G_L$ is solvable and has only the prime
divisors two and three.  It kills $C$ by
perfectness, then kills the normal group
$K/C$ of order $25$.  Every common quotient
therefore factors through the specified
$H_L$.  This proves the first fibre product
in \textup{(E.38.2)} by Goursat's lemma.
The corrected source map surjects onto
$D_{23}$ by E.24 and onto $G_L$ because
its generator images agree with an actual
Galois automorphism followed by $\rho_L$.
Its joint image lies in the stated fibre
product and has both surjective projections;
the same maximal-common-quotient argument
proves its full surjectivity.

For the four-strand comparison, first kill
the explicit perfect normal subgroup
$L^{19}$ in $T_L$.  The remaining group
$\F_2^2\rtimes A_L$, where $A_L$ is the
Nielsen image, has a faithful $251$-point
representation.  Its successive derived
indices are $1,2,6,96,96$.  The terminal
subgroup is perfect; its inverse image in
$T_L$ is perfect as an extension of perfect
groups.  Thus the maximal solvable quotient
$S_L$ has order $96$.
Its full right Cayley table and literal
representative braid words are retained.
Exhaustively traversing all $1536$ elements
of the actual binary degree-two target,
and checking all three generator edges
at each element, proves the surjection
\[
                   T_{2,2}\twoheadrightarrow S_L
 \tag{E.38.15}
\]
with the specified values of $a,b,c$.
Since $T_{23}$ is solvable and maps to
this quotient through $T_{2,2}$, the second
fibre product in \textup{(E.38.2)} is the
actual joint four-strand target.  The old
corrected triple generates $T_{23}$; the
new triple generates $T_L$ by its equality
with a Galois image.  The maximal common
quotient is $S_L$, so these triples jointly
generate the whole fibre product.  Its
kernel in $B_4$ is normal and lies in
$PB_4$, since the old binary target
already retains the standard $S_4$ quotient.

The executed files have prefix
\texttt{de\_ferrari\_sl2\_5\_} and suffixes
\texttt{model}, \texttt{nielsen}, \texttt{source},
\texttt{N}, \texttt{derived}, \texttt{coinvariants},
\texttt{modular}, and \texttt{solvable}; the
two lifted-presentation modules and their
exhaustive control are retained as well.
Every image and whole correction group
used above is a complete finite calculation.
The particular correction in
\textup{(E.38.14)} is supplied by the
proved arithmetic and product-image
argument; no uncomputed literal Galois
word is claimed as a saved certificate.
These comparisons establish the stated
finite refinements.  They do not provide
nonempty fibres over every finite
profinite quotient.
\end{proof}
